\lesson{6}{Sep 21 2021 Tue (09:21:04)}{Pure Substances and Mixtures}{Unit 1}

\subsubsection*{Classifing matter as a Mixture or Pure Substance based on its Properties}

All matter can be classified into two main categories

\begin{itemize}
    \item Mixtures
    \item Pure Substances
\end{itemize}

\begin{definition}[Matter]
    \bf{Matter} is anything that takes up \bf{Mass} and \bf{Space}.
\end{definition}

\subsubsection*{Mixtures}

A mixture is created when solids, liquids, or gases mix with one another.
Mixtures can occur with substances that are in the same states of matter or of
different states of matter. Each substance in a mixture maintains its
individual composition and property.

\begin{definition}[Homogeneous]
    In this type of mixture, just one substance looks the same throughout.
    Dissolving sugar in a glass of water forms a \bf{homogeneous} mixture,
    as the sugar seems to disappear. Even when examined under a microscope, the
    sugar water appears the same throughout.
    
    Some examples of \bf{Homogeneous Mixtures} are:
    
    \begin{itemize}
        \item Air
        \item Perfume
        \item Bleach
        \item Steel
    \end{itemize}
\end{definition}

\begin{definition}[Heterogeneous]
    Candies in a mixture will not be dispersed evenly. Heavier pieces may sink
    to the bottom of the pile, and smaller pieces may be at the edges.
    Vegetable soup is an example of a heterogeneous mixture because each
    spoonful will have a slightly different mixture of vegetable pieces and
    seasoning.
    
    Some examples of \bf{Heterogeneous Mixtures} are:
    
    \begin{itemize}
        \item Rocks
        \item Cereal
        \item Sand
    \end{itemize}
\end{definition}

\begin{definition}[Physical Separation]
    Mixtures contain compounds and elements that can be physically separated
    from one another. For example, you can pull out your favorite fruit from a
    fruit salad and leave the others in the bowl. Or you can boil a saltwater
    mixture to cause the water to evaporate and leave salt particles.
\end{definition}

\subsubsection*{Pure Substances}

A pure substance consists of a single element or type of compound. For example,
gold is a \bf{Pure Substance} because it's made from a single element:
gold.

\begin{definition}[Elements]
    An \bf{Element} is a substance that cannot be broken down into simpler
    substances. All elements on the periodic table of elements are examples of
    \bf{Pure Substances}.
\end{definition}

\begin{definition}[Compounds]
    A \bf{Compound} is matter composed of two or more \bf{Elements}. A
    molecule is the smallest unit of a \bf{Compound}. 

    Some examples of \bf{Compounds} are:
    
    \begin{itemize}
        \item Alcohol
        \item Water
        \item Baking Soda
        \item Glucose
    \end{itemize}
\end{definition}

\begin{definition}[Chemical Separation]
    Pure substances made from compounds can only be chemically separated into
    the elements that make them up. The compound table salt ($NaCl$) is made
    from the elements \bf{Sodium} and \bf{Chlorine}. A chemical
    reaction would be needed to separate these two elements from one another.
\end{definition}

\newpage
